\section{Qué se acumula antes del corte y qué queda en la frontera}\label{sec:signedmap}
Elijamos una solución uniforme y revelemos sus reinas en un orden uniforme de las filas,
independiente de la solución. Cuando quedan \(m\) filas, el tablero residual \(B_m\) registra las
casillas disponibles. En su máscara binaria \(M_m\) cabe un problema más sencillo: colocar una
torre por fila y columna. Su conteo, \(\operatorname{per}M_m\), ya incorpora las restricciones
que imponen las piezas reveladas. Para obtener \(Q(B_m)\) falta exigir que las nuevas piezas
tampoco se ataquen entre sí por las diagonales.

Esto sugiere una primera decisión. Sean \(L\) y \(L_T\) los logaritmos de los predictores de reinas
y de torres en el tablero residual. En la parte de torres podemos usar el conteo exacto:
\[
 g_T(M)=\log\operatorname{per}M-L_T(M),\qquad L^\sharp(B)=L(B)+g_T(M).
\]
El ajuste se hace sobre la misma máscara en cada paso. Así una diferencia entre dos tableros hijos
no mezcla el error de contar torres con el coste de prohibir un ataque diagonal. En el tablero
inicial \(n g_T(M_n)\) tiende a \(F(1)-1/12\). La máscara llena es separable y \(F(1)=0\); de
aquí procede el primer sumando de~\eqref{eq:main}.

\paragraph{La suma hasta la parada.}
Fijamos los márgenes de regularidad y el parámetro macroscópico \(\delta\), y después un
\(b>0\) suficientemente pequeño. Paramos en la primera salida de la clase regular o al alcanzar
\(j=\lfloor bn\rfloor\), como en la comparación sobre historias detenidas de~\cite{tc}. Sea \(F_b\)
la llegada regular y completable a ese nivel.
En cada fila, la ley adaptativa elige una casilla con probabilidad proporcional al peso de su
tablero hijo. Si la suma de esos pesos coincidiera con el peso del padre, la recursión se
cumpliría exactamente en esa fila. El logaritmo de su cociente mide el defecto de ese paso.
Llamamos \(A^\sharp\) a la suma de estos logaritmos hasta la parada.

La solución uniforme y la ley adaptativa utilizan la misma regla de parada. Esto permite comparar
las probabilidades de una misma historia y conservar el signo de la diferencia de las salidas.
El Apéndice~\ref{app:signed} da la identidad exacta y el Apéndice~\ref{app:arrival} controla
el resto en media. A \(b\) fijo, el defecto acumulado en el interior aporta
\(\int_b^1 k(x)x^{-2}\,dx\). El Apéndice~\ref{app:interior} identifica sus dieciséis términos.

\paragraph{Lo que sobrevive en el tablero de llegada.}
Consideremos dos casillas disponibles \(e,f\) en una misma diagonal. Están en filas y columnas
distintas, de modo que dos torres podrían ocuparlas a la vez. Las otras dos posiciones del
cuadrado determinan cuatro estados. Para \(e=(r,q)\) y \(f=(r+h,q+h)\), con \(h>0\),
sean \(X\) e \(Y\) los indicadores de disponibilidad de \((r,q+h)\) y \((r+h,q)\),
respectivamente. La Figura~\ref{fig:frontera} muestra los cuatro casos.
\begin{figure}[htbp]
\centering
\begin{tikzpicture}[x=.58cm,y=.58cm]
 \foreach \dx/\sx/\sy in {0/0/0,5.8/1/0,11.6/0/1,17.4/1/1}{
  \begin{scope}[shift={(\dx,0)}]
   \bcuadricula{4}
   \draw[btinta,line width=.9pt] (1.5,3.5)--(3.5,1.5);
   \ifnum\numexpr\sx+\sy\relax=2\relax
    \draw[btinta,dash pattern=on 2pt off 1.5pt,line width=.9pt] (1.5,1.5)--(3.5,3.5);
   \fi
   \foreach \fr/\fc/\libre in {1/2/1,3/4/1,1/4/\sx,3/2/\sy}{
    \begin{scope}[shift={(\fc-.5,4.5-\fr)}]
     \fill[bpapel] (0,0) circle[radius=.26];
     \ifnum\libre=1\relax
      \draw[btinta,line width=1pt] (0,0) circle[radius=.22];
     \else
      \draw[btinta,line width=1.05pt] (-.17,-.17)--(.17,.17) (-.17,.17)--(.17,-.17);
     \fi
    \end{scope}
    \ifnum\libre=1\relax \bmarca{4}{\fr}{\fc}\fi
   }
   \node[anchor=north,font=\small] at (2,-.18) {$(X,Y)=(\sx,\sy)$};
  \end{scope}
 }
\end{tikzpicture}
\caption{Los cuatro estados del cuadrado de frontera se dibujan en una ventana de
\(4\times4\) de un tablero mayor, con coordenadas locales \(r=1\), \(q=2\), \(h=2\).
El círculo indica una casilla disponible y la cruz una vetada. La línea continua une
\(e=(1,2)\) y \(f=(3,4)\). La discontinua une \((1,4)\) con \((3,2)\) cuando ambas
están disponibles. Las demás casillas se muestran sólo como referencia. En el estado \(11\),
tras restar los efectos de cada par, los coeficientes de interacción en el logaritmo del conteo
y en el del predictor son \(-1\) y \(+2\), en unidades del producto de los cuatro pesos al orden
principal. Su diferencia \(-1-(+2)=-3\) da el signo de \(-S\) en~\eqref{eq:exceso-cuadrado}.}
\label{fig:frontera}
\end{figure}


Los cuatro casos se realizan en la solución \((1,3,5,7,2,4,6)\) del tablero de orden siete.
Antes de revelar una reina, \(X=Y=1\). Revelar \((7,6)\) veta sólo \(Y\) entre los cuatro
vértices y revelar \((4,7)\) veta sólo \(X\). Si se revelan ambas, siguen disponibles \(e,f\).

La respuesta del par marcado se evalúa con \(X,Y\) fijados. En el estado \(11\) se puede
marcar cualquiera de los dos pares diagonales, y cada uno aporta su propia respuesta. Ambos
pertenecen al mismo cuadrado, por lo que su interacción entra una sola vez en la suma sobre
cuadrados completos. Denotamos por \(B_n(b)\) la suma de las tres medias de frontera.
Al tomar \(n\to\infty\) y después \(b\downarrow0\), las dos medias sobre pares dan \(P_T\)
y \(J\), y la media de interacción da \(-S\). Los Apéndices~\ref{app:boundary}
y~\ref{app:weights} separan las contribuciones. El Apéndice~\ref{app:transfer}
efectúa el promedio de la respuesta del par bajo la ley de llegada.

\paragraph{Por qué se comparan dos cortes.}
Para \(H=P_T+J-S\), escribamos
\[
 \Phi(b)=F(1)-\tfrac1{12}+\int_b^1 k(x)x^{-2}\,dx+H.
\]
El error cumple \(\limsup_n |nR_Q(n)-\Phi(b)|\le C_\delta b+\eta(b)\), donde
\(\eta(b)=\limsup_n|B_n(b)-H|\to0\) al bajar \(b\). La constante \(C_\delta\) sirve para
todos los cortes admisibles suficientemente pequeños. Cada elección de \(b\) aproxima, por tanto,
la misma sucesión \(nR_Q(n)\). Al comparar \(b\) y \(b'\), la diferencia entre las aproximaciones
tiende a cero con los dos cortes. Así se obtiene un límite común antes de integrar hasta el
extremo. Por analiticidad y \(k(0)=0\), un término lineal de \(k\) produciría allí la divergencia
\(k'(0)\log(1/b)\). La convergencia obliga a que desaparezca. El Apéndice~\ref{app:closing}
detalla la comparación y deduce el corolario para \(d'(0)\).

La Figura~\ref{fig:cortes} ilustra los dos niveles en una historia concreta. En la prueba,
cada corte admisible da una aproximación mediante promedios sobre historias. Como ambas
aproximaciones se refieren al mismo \(nR_Q(n)\), la desigualdad triangular da
\(\lvert\Phi(b)-\Phi(b')\rvert\le C_\delta(b+b')+\eta(b)+\eta(b')\).

\begin{figure}[htbp]
\centering
\begin{tikzpicture}[x=.58cm,y=.58cm]
 \foreach \etapa/\dx in {2/0,4/10.4}{
  \begin{scope}[shift={(\dx,0)}]
   \bcuadricula{6}
   \breina{6}{1}{2}
   \breina{6}{2}{4}
   \ifnum\etapa=2\relax
    \def\casillas{3/1,3/6,4/1,4/3,5/3,5/5,6/1,6/3,6/5,6/6}
    \def\corte{b=2/3}
   \else
    \breina{6}{3}{6}
    \breina{6}{4}{1}
    \def\casillas{5/3,5/5,6/5}
    \def\corte{b'=1/3}
   \fi
   \foreach \fila/\col in \casillas{
    \draw[btinta,line width=.55pt] (\col-.5,6.5-\fila) circle[radius=.16];
    \bmarca{6}{\fila}{\col}
   }
   \draw[btinta,line width=.95pt] (0,6-\etapa)--(6,6-\etapa);
   \node[font=\small,anchor=north] at (3,-.18) {$\corte,\quad m=\the\numexpr6-\etapa\relax$};
  \end{scope}
 }
 \draw[->,>=stealth,btinta,line width=.8pt] (6.55,3)--(9.45,3)
   node[midway,above,font=\scriptsize] {filas 3 y 4};
\end{tikzpicture}
\caption{Dos cortes de la solución \((2,4,6,1,3,5)\). La entrada \(r\) da la columna de la reina
en la fila \(r\). En el ejemplo, las filas se revelan de arriba abajo y la raya horizontal marca hasta dónde
se ha llegado. Las coronas indican las reinas reveladas y los círculos con punto, las casillas
que esas reinas dejan disponibles. Al revelar las filas \(3\) y \(4\), quedan tres de las diez.}
\label{fig:cortes}
\end{figure}

